\section*{Executive Summary}
\addcontentsline{toc}{section}{Executive Summary}

When a fault occurs in a substation, protective relays energize breaker trip coils,
and reclosing and motor-operated equipment may follow. All of this equipment draws
current from the station's DC battery system. Engineers must find the largest
\emph{momentary} current each DC supply may have to deliver. Today, that means
tracing protection zones on paper drawings and adding up coil currents in a
spreadsheet.

\app{} is a desktop application that does this work on a rebuilt copy of the
one-line diagram. The user redraws the station over the original PDF and enters each
device's coil and motor currents. The tool then:

\begin{itemize}
  \item proposes bus, transformer and line protection zones by tracing the wiring
        between current transformers (CTs), taking open switches into account;
  \item lets the engineer review each zone, assign the equipment that operates at
        each stage (primary trip, reclose, motor inrush and so on), and verify it;
  \item adds up the currents that flow at the same time for each zone, stage and DC
        supply, and reports the peak current each supply must deliver;
  \item exports a PDF calculation report with an annotated one-line diagram.
\end{itemize}

Results stay marked \emph{provisional} until every zone has been verified, so an
unfinished study cannot be mistaken for a final one. The tool calculates momentary
current only; it does not size battery capacity in ampere-hours.

\app{} is written in Python with the Qt user-interface toolkit and is delivered as a
portable Windows package that needs no installation. Its results matched an
independent hand calculation exactly on all 13 provided layouts, and 288 automated
tests run on Windows and Linux whenever the code changes. The sponsor reviewed the
provided layouts and confirmed that they are correct.
